\section*{Разбор задачи из тестового варианта №5}
\addcontentsline{toc}{section}{Разбор задачи из тестового варианта №5}
\begin{Task}{Условная плотность и вынесение известного сомножителя}{}
    Найдите условную плотность $X$ относительно $Y/X$, если случайные величины $X$ и $Y$ независимы, $X \sim \text{U}(0, 1)$, а $Y$ имеет плотность
    \[
        p_Y(x) = 4x^3 \cdot \I\{x \in (0, 1)\}.
    \]
    Используя полученную условную плотность, вычислите условное математическое ожидание $\E(XY|Y/X)$.
\end{Task}

\begin{Solution}
    \textbf{Семантический анализ задачи.} Вычисление условных распределений абсолютно непрерывных векторов требует перехода к новой системе координат через якобиан преобразования. Фундаментальной сложностью задачи является нелинейность преобразования (отношение величин), что приводит к зависимости границ носителя одной переменной от значения другой. Для оптимизации расчета математического ожидания необходимо использовать свойство измеримости (вынесение известного условия за знак ожидания).

    \textbf{Шаг 1. Совместная плотность исходного вектора.}
    Поскольку случайные величины $X$ и $Y$ независимы, их совместная плотность равна произведению маргинальных плотностей. С учетом равномерного распределения $X \sim \text{U}(0, 1)$, где $p_X(x) = 1 \cdot \I\{x \in (0, 1)\}$, получаем:
    \[
        p_{X,Y}(x, y) = p_X(x)p_Y(y) = 4y^3 \cdot \I\{x \in (0, 1), y \in (0, 1)\}.
    \]

    \textbf{Шаг 2. Переход к новым переменным и вычисление якобиана.}
    Введем новую случайную величину $Z = Y/X$. Для использования теоремы о замене переменных дополним ее до двумерного преобразования, оставив $X$ без изменений. Выразим старые переменные через новые:
    \[
        \begin{cases}
            X = x \\
            Y = xz
        \end{cases}
    \]
    Вычислим якобиан данного преобразования:
    \[
        J = \det \begin{pmatrix}
            \frac{\partial x}{\partial x} & \frac{\partial x}{\partial z} \\[1em]
            \frac{\partial (xz)}{\partial x} & \frac{\partial (xz)}{\partial z}
        \end{pmatrix} 
        = \det \begin{pmatrix}
            1 & 0 \\
            z & x
        \end{pmatrix} = x.
    \]
    Совместная плотность новой пары $(X, Z)$ имеет вид $p_{X,Z}(x, z) = p_{X,Y}(x, xz) \cdot |J|$. Поскольку по условию $x \in (0, 1)$, модуль раскрывается положительно ($|x| = x$):
    \[
        p_{X,Z}(x, z) = 4(xz)^3 \cdot x = 4x^4z^3.
    \]

    \textbf{Шаг 3. Анализ носителя (области ненулевой плотности).}
    Определим границы, в которых совместная плотность $p_{X,Z}(x, z)$ отлична от нуля. Подставим замену $y = xz$ в исходные ограничения $x \in (0, 1)$ и $y \in (0, 1)$:
    \[
        \begin{cases}
            0 < x < 1 \\
            0 < xz < 1
        \end{cases} \implies
        \begin{cases}
            0 < x < 1 \\
            x < \frac{1}{z} \quad (\text{так как } z > 0)
        \end{cases}
    \]
    Отсюда следует принципиально важный вывод: верхняя граница для $x$ зависит от значения $z$ и равна $\min\left(1, \frac{1}{z}\right)$. Пространство значений $z$ разбивается на два непересекающихся интервала:
    \begin{itemize}
        \item При $z \in (0, 1]$ граница интегрирования по $x$ определяется условием $x \in (0, 1)$.
        \item При $z \in (1, +\infty)$ граница интегрирования по $x$ сужается до $x \in \left(0, \frac{1}{z}\right)$.
    \end{itemize}

    \textbf{Шаг 4. Вычисление маргинальной и условной плотностей.}
    Найдем маргинальную плотность $p_Z(z)$, интегрируя совместную плотность по $x$ с учетом найденных границ.

    \textit{Случай 1: $z \in (0, 1]$.}
    \[
        p_Z(z) = \int_{0}^{1} 4x^4z^3 \, dx = 4z^3 \left[ \frac{x^5}{5} \right]_0^1 = \frac{4}{5}z^3.
    \]
    Условная плотность на этом интервале:
    \[
        p_{X|Z}(x|z) = \frac{p_{X,Z}(x, z)}{p_Z(z)} = \frac{4x^4z^3}{\frac{4}{5}z^3} = 5x^4, \quad x \in (0, 1).
    \]

    \textit{Случай 2: $z \in (1, +\infty)$.}
    \[
        p_Z(z) = \int_{0}^{1/z} 4x^4z^3 \, dx = 4z^3 \left[ \frac{x^5}{5} \right]_0^{1/z} = 4z^3 \cdot \frac{1}{5z^5} = \frac{4}{5z^2}.
    \]
    Условная плотность на этом интервале:
    \[
        p_{X|Z}(x|z) = \frac{p_{X,Z}(x, z)}{p_Z(z)} = \frac{4x^4z^3}{\frac{4}{5z^2}} = 5x^4z^5, \quad x \in \left(0, \frac{1}{z}\right).
    \]

    \textbf{Шаг 5. Вычисление условного математического ожидания.}
    Требуется найти $\E(XY | Z)$. Выразим аргумент функции через переменные $X$ и $Z$: $XY = X(XZ) = X^2Z$. 
    
    Воспользуемся свойством измеримости: относительно условия $Z$, величина $Z$ является константой и может быть вынесена за знак условного математического ожидания.
    \[
        \E(XY | Z) = \E(X^2Z | Z) = Z \cdot \E(X^2 | Z).
    \]
    Сведение задачи к вычислению $\E(X^2 | Z)$ существенно упрощает интегрирование. Рассчитаем ожидание для каждого интервала значений $Z$.

    \textit{Для $Z \in (0, 1]$:}
    \begin{align*}
        \E(X^2 | Z=z) &= \int_{0}^{1} x^2 \cdot (5x^4) \, dx = 5 \int_{0}^{1} x^6 \, dx = \frac{5}{7}. \\
        \E(XY | Z=z) &= z \cdot \frac{5}{7} = \frac{5}{7}z.
    \end{align*}

    \textit{Для $Z \in (1, +\infty)$:}
    \begin{align*}
        \E(X^2 | Z=z) &= \int_{0}^{1/z} x^2 \cdot (5x^4z^5) \, dx = 5z^5 \int_{0}^{1/z} x^6 \, dx = 5z^5 \left[ \frac{x^7}{7} \right]_0^{1/z} = \frac{5z^5}{7z^7} = \frac{5}{7z^2}. \\
        \E(XY | Z=z) &= z \cdot \frac{5}{7z^2} = \frac{5}{7z}.
    \end{align*}

    Обобщая полученные результаты, записываем итоговую функцию от случайной величины $Z = Y/X$.
\end{Solution}

\begin{Answer}
    Условная плотность $X$ при условии $Z = Y/X$:
    \[
        p_{X|Z}(x|z) = \begin{cases}
            5x^4 \cdot \I\{x \in (0, 1)\}, & 0 < z \le 1, \\[0.5em]
            5x^4z^5 \cdot \I\left\{x \in \left(0, \frac{1}{z}\right)\right\}, & z > 1.
        \end{cases}
    \]
    Условное математическое ожидание:
    \[
        \E(XY | Y/X) = \begin{cases}
            \frac{5}{7}\left(\frac{Y}{X}\right), & 0 < \frac{Y}{X} \le 1, \\[1em]
            \frac{5}{7\left(\frac{Y}{X}\right)}, & \frac{Y}{X} > 1.
        \end{cases}
    \]
\end{Answer}
